We studied the de la Vega--Lueker and Karmarkar--Karp algorithms,
explaining their grouping and rounding techniques and deriving their
approximation guarantees and running-time bounds. The de la Vega--Lueker
algorithm gives an APTAS using at most
$(1+\eps)\OPT(I)+O(\eps^{-2})$ bins for any fixed $\eps>0$.
Karmarkar--Karp achieves the stronger additive guarantee
$\OPT(I)+O(\log^2(\OPT(I)+1))$ in polynomial time.

By studying these two algorithms, we learned how grouping reduces the number of distinct item sizes, allowing a polynomial-size configuration ILP for fixed $\eps$ in the
de la Vega--Lueker scheme. LP relaxation and careful rounding then
produce a feasible packing with a provable approximation guarantee.

Beyond our scope, Hoberg and Rothvoss~\cite{hobergLogarithmicAdditive2017}
improved the additive guarantee to $O(\log\OPT(I))$ using discrepancy
theory and a two-stage construction that packs items into containers
and then containers into bins. This remains the best known polynomial-time
additive approximation guarantee.
Understanding these discrepancy methods would be a natural extension
of our study of LP rounding.

Recent work also establishes limits on further improvements.
OpenAI's September 2026 preprint~\cite{openaiAdditiveHardness2026}
constructs instances with arbitrarily large gaps
$\OPT(I)-\LIN(I)$, disproving the Modified Integer Round-Up
Conjecture $\OPT(I)\le\lceil\LIN(I)\rceil+1$.
It also proves that no deterministic polynomial-time algorithm can
guarantee $\OPT(I)+C$ bins for any fixed constant $C$, unless
$\mathrm{P}=\mathrm{NP}$; both results hold even when every item
exceeds $1/6$. These results show that a constant additive loss is
unattainable in general under $\mathrm{P}\ne\mathrm{NP}$, while leaving
room to improve the logarithmic upper bound. Their proofs are beyond
the scope of this report.
